\section{Data, Corruption Protocol and Metrics}\label{sec:data}
\textbf{Oxford-IIIT Pet (Tasks 1--3).} The official train/validation collection (3,680 images~\cite{parkhi2012}) was split 80/20 with seed 42 into \textbf{2,944 training and 736 validation images}; the 3,669 official test images were never used for training or tuning. Labels are ignored, images are converted to RGB and resized (bicubic) to $128\times128$; one split file serves all three tasks.

\textbf{Runtime corruption.} Nothing corrupted is stored. Each time a training image is loaded, one of four conditions (clean, salt-and-pepper, blur, occlusion) is drawn with equal probability and a fresh specification is sampled (Table~\ref{tab:corruption}). The same function applies it in training, validation and testing; the classifier uses exactly balanced batches ($B/4$ per class).

\begin{table*}[t]
\centering
\caption{Complete corruption configuration. Training values are re-sampled at every load; validation uses the random ranges once with fixed seeds; testing uses three fixed severities.}\label{tab:corruption}
\small
\begin{tabular}{p{2.5cm}p{6.2cm}p{7.4cm}}
\toprule
Condition & Training / validation (random) & Fixed test severities (low, medium, high) \\
\midrule
Clean & original image & original image \\
Salt-and-pepper & $p\sim U(0.02,0.15)$; exactly $\mathrm{round}(pHW)$ pixels (all channels) set to black or white with equal probability & $p=0.03,\;0.08,\;0.15$ \\
Gaussian blur & kernel $\{3,5,7\}$, $\sigma\sim U(0.5,2.5)$, separable Gaussian, reflect padding & $(k,\sigma)=(3,0.7),\;(5,1.5),\;(7,2.5)$ \\
Rectangular occlusion & 1--3 black rectangles, union area 10--35\% of the image, random positions and aspect ratios & 1, 2, 3 rectangles covering 10\%, 20\%, 35\% ($\pm1\%$) \\
\bottomrule
\end{tabular}
\end{table*}

\begin{figure*}[t]
\centering
\includegraphics[width=0.8\textwidth]{corruption_grid.pdf}
\caption{The ten deterministic test conditions for one test image (clean and three severities of each corruption).}\label{fig:corr}
\end{figure*}

\textbf{Manifests.} Validation (736 entries, 184 per class) and test manifests are generated once and store the corruption type, severity, parameters, mask coordinates and seed, with a SHA-256 digest. Every test image appears under all ten conditions, giving \textbf{36,690 test inputs} (3,669 images $\times$ 10); all test tables use this full manifest.

\textbf{FS2K (Task 4).} The official training/test definitions are used. 15\% of the training portion is held out as validation, stratified by style with seed 42, leaving \textbf{899 training and 159 validation pairs}; the \textbf{1,046 official test pairs} (619, 381 and 46 in styles 1, 2, 3) are used only for the final evaluation. Photographs are $128\times128$ RGB, sketches $128\times128$ grayscale; augmentation (flip, random crop-and-resize) is applied identically to both members of a pair.

\textbf{Metrics.} L1 (images in $[0,1]$), SSIM (11-pixel Gaussian window) and PSNR in dB, clipped at 60\,dB because identical images give infinite PSNR. The score of the untouched corrupted input (``no restoration'') is always shown, since a restoration model is useful only where it beats it. Classification uses accuracy, macro precision/recall/F1, per-class values and a row-normalised confusion matrix.
